\documentclass[xcolor={dvipsnames},handout]{beamer}

\usepackage{assets/pres-template_MOW}
\usepackage{pdfpages}

%--------------------------------------------------------------------------
% Specific to this document ---------------------------------------
%--------------------------------------------------------------------------
% \setbeamercovered{transparent}

%%%%%%%%%%%%%%%%%%%%%%%%%%%%%%%%%%%
\setlength{\tabcolsep}{1.3pt}
\title{PLSC 48110/28110}
\subtitle{Lecture 3: Behavioral games}
\date{}
\author{Molly Offer-Westort}
\institute{Department of Political Science, \\University of Chicago}


\usepackage{assets/Sweave}
\hypersetup{pdftitle={PLSC 48110/28110 - Lecture 3: Behavioral games},pdfauthor={Molly Offer-Westort},pdfpagemode=UseNone}
\begin{document}

%-------------------------------------------------------------------------------%
\frame{\titlepage
\thispagestyle{empty}
}
%-------------------------------------------------------------------------------%
\begin{frame}{Behavioral/laboratory games}


\begin{itemize}
\item What is the laboratory?
\item Lab games as an approach to measurement
\begin{itemize}
\item Context-free measures of "economic" preferences 
\item "Empirical anthropology" 
\item Types of things that may be measured
\begin{itemize}
\item Risk preferences: risk aversion/neutrality/loving 
\item Time preferences: patience, discounting 
\item Social preferences: altruism, pro-social behavior, trust
\item Cognitive skills: behavior under cognitive load
\end{itemize}
\end{itemize}
\end{itemize}

\end{frame}


%%%%%NOTE%%%%%
\note{
\scriptsize \singlespacing

\begin{itemize}
\item xxxx
\end{itemize}

}

%-------------------------------------------------------------------------------%
\begin{frame}{Benefits of this approach to measurement}


\begin{itemize}
\item Closer to real behavior (?) than survey experiments
\item Fully controlled environment; removes other contextual information, social norms 
\item Useful for measuring heterogeneity 
\item Survey fatigue? Respondent engagement?
\end{itemize}

\end{frame}


%%%%%NOTE%%%%%
\note{
\scriptsize \singlespacing

\begin{itemize}
\item xxxx
\end{itemize}

}
%-------------------------------------------------------------------------------%
\begin{frame}{Terminology}

\begin{itemize}
\item "rational-actor" framework; 
\begin{itemize}
\item "the action chosen by a decision-maker is at least as good, according to her preferences, as every other available action" \cite[p.6]{osborne2004introduction}. 
\end{itemize}
\item "rational" preferences must be: 
\begin{itemize}
\item Complete;  for any two actions, a person knows whether she prefers the outcome under action $a$ to action $b$
\item Transitive;  if a person prefers the outcome under action $a$ to that under action $b$, and $b$ to $c$, she must prefer $a$ to $c$
\item Ordinal;  only the relative, not the absolute, value of utility matters (doesn't matter \emph{how much} more she prefers outcome under action $a$ over $b$)
\end{itemize}
\item What does it mean for players to be "self-interested?"
\item Does rational action require self-interest?
\end{itemize}

\end{frame}


%%%%%NOTE%%%%%
\note{
\scriptsize \singlespacing

\begin{itemize}
\item xxxx
\end{itemize}

}

%-------------------------------------------------------------------------------%
\begin{frame}{Terminology}

\begin{itemize}
\item One-shot vs. repeated play 
\item Subgame perfect equilibrium 
\item Strategy method vs. dynamic interactions
\end{itemize}

\end{frame}


%%%%%NOTE%%%%%
\note{
\scriptsize \singlespacing

\begin{itemize}
\item xxxx
\end{itemize}

}

%-------------------------------------------------------------------------------%
\begin{frame}{Some specific games}

From \cite{camerer2004measuring}

\end{frame}


%%%%%NOTE%%%%%
\note{
\scriptsize \singlespacing

\begin{itemize}
\item xxxx
\end{itemize}

}

%-------------------------------------------------------------------------------%

{
\setbeamercolor{background canvas}{bg=}
\includepdf[pages=8-9]{assets/fehr_2004.pdf}
}
%-------------------------------------------------------------------------------%
%%% [[TEMPLATE]] %%%
\begin{frame}{Variations}

\begin{itemize}
\item Ultimatum game w/ competition
\item Opportunities for punishment 
\item Public goods game with partners vs. strangers
\end{itemize}

\end{frame}


%%%%%NOTE%%%%%
\note{
\scriptsize \singlespacing

\begin{itemize}
\item xxxx
\end{itemize}

}

%-------------------------------------------------------------------------------%

\begin{frame}{Ethical considerations}

\begin{itemize}
\item Pay compared to time spent  (1-2 days' wages...several months' wages?)
\item Social consequences
\item Other concerns?
\end{itemize}

\end{frame}


%%%%%NOTE%%%%%
\note{
\scriptsize \singlespacing

\begin{itemize}
\item xxxx
\end{itemize}

}

%-------------------------------------------------------------------------------%
\begin{frame}{Relevance}

\begin{itemize}
\item Why is any of this relevant for political science/policy? 
\item How do you know you're measuring the "right" thing, and that respondents are interpreting the game the way you want them to?
\item Can we compare across societies?
\end{itemize}

\end{frame}


%%%%%NOTE%%%%%
\note{
\scriptsize \singlespacing

\begin{itemize}
\item xxxx
\end{itemize}

}

%-------------------------------------------------------------------------------%
\backupbegin
%-------------------------------------------------------------------------------%

\begin{frame}[allowframebreaks]{References}
    \bibliographystyle{apalike}
    \bibliography{assets/bib}
\end{frame}
%-------------------------------------------------------------------------------%
\backupend
\end{document}
